\section{The hyperk\"ahler endpoint}\label{sec:endpoint}

We now use the auxiliary form $D$ to construct the endpoint while the
actual triples retain the three original classes.  The first K\"ahler
step produces $C_1\in[C]$; the uniform chamber estimate keeps $[B]$
available along the segment from $C$ to $C_1$; and the second K\"ahler
step produces its prescribed endpoint $B_1\in[B]$.  We then choose
the third form smoothly along that segment with both endpoint values
fixed, using \Cref{lem:smooth-tamers}.  We first record the classical
class and volume criteria used in the two K\"ahler steps.

\begin{lemma}\label{lem:endpoint-kahler-class}
Let $(X,J)$ be a compact connected complex surface with even first
Betti number. Suppose a real $(1,1)$ class $H$ has $H^2>0$ and contains
a closed form $\beta$ taming $J$. Then $H$ is a K\"ahler class.
If $\nu$ is a smooth positive four-form with $\int_X\nu=H^2$, there is
a K\"ahler form $\beta_1\in H$ with $\beta_1^2=\nu$.
\end{lemma}

\begin{proof}
The compact-surface K\"ahler criterion gives a K\"ahler form $\kappa$
because $b_1$ is even \cite[Theorem~11]{Buchdahl1999}; see also
\cite[Corollaire~5.7]{Lamari1999}. Taming gives
\[
 H\cdot[D]>0
 \quad\text{for every nonzero effective curve }D,
 \qquad H\cdot[\kappa]>0.
\]
For the second inequality, pointwise
$\beta\kappa=\beta^{1,1}\kappa>0$, since the invariant part of a tamer
is positive definite. Together with $H^2>0$, these inequalities are
the numerical K\"ahler-class criterion on a surface
\cite[Corollary~15]{Buchdahl1999}. Equivalently, the segment from
$[\kappa]$ to $H$ has positive square and remains positive on every
curve, so the component condition in the numerical characterization
of the K\"ahler cone is satisfied
\cite[Theorem~0.1]{DemaillyPaun2004}.

Choose a K\"ahler representative $\kappa_H\in H$ and set
$f=\log(\nu/\kappa_H^2)$. The volume hypothesis says
$\int_X e^f\kappa_H^2=\int_X\kappa_H^2$. Yau's theorem produces a
smooth cohomologous K\"ahler form $\beta_1$ with
$\beta_1^2=e^f\kappa_H^2=\nu$
\cite[Section~4, Theorem~1]{Yau1978}.
\end{proof}

\begin{lemma}\label{lem:endpoint-hodge-type}
Let $X$ be a compact connected K\"ahler surface with a nowhere-zero
holomorphic two-form $\Omega$. A real degree-two class $H$ is of type
$(1,1)$ if and only if it is intersection-orthogonal to
$[\Re\Omega]$ and $[\Im\Omega]$.
\end{lemma}

\begin{proof}
Every holomorphic two-form is $f\Omega$ for a global holomorphic
function $f$, hence is a constant multiple of $\Omega$. Thus
$H^{2,0}(X)=\C[\Omega]$. The real plane generated by its real and
imaginary parts is nondegenerate and positive for the intersection
form: their squares are equal and positive and their mixed product
vanishes. Hodge decomposition and type considerations make its
orthogonal complement exactly $H^{1,1}(X;\R)$.
\end{proof}

\begin{figure}[htbp]
\centering
\begin{tikzpicture}[
  box/.style={draw=blue!45!black,rounded corners=2pt,
    fill=blue!3,align=center,text width=3.35cm,minimum height=1.05cm,
    inner sep=5pt,font=\small},
  aux/.style={box,draw=green!35!black,fill=green!3,text width=4.1cm},
  route/.style={-{Stealth[length=2mm]},thick,draw=blue!50!black},
  aid/.style={-{Stealth[length=2mm]},dashed,draw=green!35!black},
  lab/.style={font=\scriptsize,align=center,fill=white,inner sep=2pt}
]
\node[box] (start) at (0,0) {Prepared triple\\$A,C,B$};
\node[box] (move) at (4.65,0) {Moving triple\\$A,C_t,B_t$};
\node[box] (finish) at (9.3,0) {Hyperk\"ahler triple\\$A,C_1,B_1$};
\draw[route] (start) -- node[lab,above=20pt] {$C_t\in[C]$\\$B_t\in[B]$} (move);
\draw[route] (move) -- node[lab,above=20pt] {prescribed\\endpoint} (finish);
\node[aux] (auxiliary) at (4.65,-2.25)
  {Auxiliary closed form $D$\\$AD=CD=0$,\quad $D^2=A^2$};
\draw[aid] (start.south) -- ++(0,-1.725) --
  node[lab,below] {small fiber area} (auxiliary.west);
\draw[aid] (auxiliary.north) -- node[lab,right]
  {K\"ahler step for $A\pm iD$\\in $[C]$; taming criterion for $[B]$} (move.south);
\draw[aid] (auxiliary.east) -- (finish.south |- auxiliary.east) --
  node[lab,right] {using $C_1$ and\\a tamer in $[B]$:\\K\"ahler step for\\$A\pm iC_1$} (finish.south);
\end{tikzpicture}
\caption{The final deformation. Solid arrows trace the deformation; dashed
arrows indicate auxiliary dependencies. The upper row consists of
actual triples in the original classes. The auxiliary form supplies the
first integrable pair and the numerical control that keeps $[B]$
available along the moving pair. After both K\"ahler steps, the tamers
$B_t$ are chosen with initial value $B$ and final value $B_1$.}
\label{fig:endpoint-path}
\end{figure}


\begin{proof}[Proof of \Cref{thm:main}]
We trace the actual form paths and then construct the endpoint;
\Cref{fig:endpoint-path} summarizes the final stages.
All coefficient changes below are constant orthogonal changes, and
will be undone at the end.

\paragraph{\emph{The prepared pair.}}
Choose the primitive null class $F$ of \Cref{prop:chamber-class} and
rotate the coefficient basis so that $[B]$ points in the direction of
$F^+$ and the pair $[A],[C]$ spans its orthogonal plane in $P$.
The rotated initial triple is denoted $(A,C,B)$; its cohomology Gram
matrix is still the identity.

The families result and \Cref{prop:nonsplitting,thm:pencil} give an
exact pair deformation to a torus fibration with the stated standard
nodal neighborhoods, keeping $B$ fixed.
\Cref{prop:regular-preparation} then gives a further path of positive
closed triples to invariant regular data with the auxiliary form
$\Psi$.  Both stages preserve the three individual classes and start
at the actual representatives supplied by the preceding stage.
Retain the notation
$(A,C,B)$ for the resulting triple. In particular
\begin{equation}\label{eq:endpoint-periods}
 [A]^2=[C]^2=[B]^2=1,\qquad
 [A][C]=[A][B]=[C][B]=0.
\end{equation}

\paragraph{\emph{The auxiliary solution.}}
Fix a sufficiently small $\epsilon>0$ for
\Cref{thm:auxiliary-solution,prop:auxiliary-class}, and set
$D=D_\epsilon$. The form $D$ is closed and smooth and satisfies
\begin{equation}\label{eq:endpoint-auxiliary}
 AD=CD=0,\qquad D^2=A^2>0.
\end{equation}
The triple with forms $A,C,D$ is positive. The same choice of
$\epsilon$ gives the numerical conclusion of
\Cref{prop:auxiliary-class} for every later definite pair with the
classes $[A],[C]$. These choices precede the K\"ahler constructions
that follow.

\paragraph{\emph{The first K\"ahler step.}}
By \Cref{lem:complex-symplectic}, $A\pm iD$ is a nowhere-zero
holomorphic two-form for an integrable complex structure $I$; choose the
sign so that $C$ tames $I$. This is possible by
\Cref{lem:definite-pair} and positivity of $(A,D,C)$.
\Cref{prop:topology} gives even $b_1$, so the surface is K\"ahler.
The period relations \eqref{eq:endpoint-periods} and
\eqref{eq:endpoint-auxiliary} give
$[C][A]=[C][D]=0$. Therefore \Cref{lem:endpoint-hodge-type} puts
$[C]$ in real type $(1,1)$.

Apply \Cref{lem:endpoint-kahler-class} with
$H=[C]$ and $\nu=A^2$. Its integral condition holds by
\eqref{eq:endpoint-periods}. We obtain a K\"ahler representative
$C_1\in[C]$ satisfying
\begin{equation}\label{eq:endpoint-first-yau}
 AC_1=DC_1=0,\qquad C_1^2=A^2.
\end{equation}
Let $C_t=(1-t)C+tC_1$. Both endpoints tame $I$, so the entire segment
does. Thus $(A,D,C_t)$ is
positive, and $(A,C_t)$ is a definite pair in the original pair
classes. Moreover $D$ tames its associated almost-complex structure
with the continuously chosen sign. Indeed $D$ is wedge-orthogonal
to both pair forms by \eqref{eq:endpoint-auxiliary} and
\eqref{eq:endpoint-first-yau}, and has positive square.

\paragraph{\emph{Keeping the third class.}}
Let $J_t$ now denote the structure associated with the moving pair
$(A,C_t)$ and tamed by $D$. For every irreducible $J_t$ curve, its
class $d$ has $[D]d>0$. \Cref{prop:auxiliary-class} implies
$[B]d>0$, and places $[B]$ and $[D]$ in the same open timelike cone
of the fixed space $V$ from \eqref{eq:pair-lorentz-space}.
The hypotheses of \Cref{thm:taming-cone} therefore hold with
$U=[D]$ and $H=[B]$. Thus $[B]$ contains a tamer for every $J_t$.

At $t=0$ this sign of $J_t$ agrees with the sign previously tamed by
the actual third form $B$.  To check this pointwise, take a regular
fiber with its prepared orientation.  Both $B$ and $D$ restrict
positively to its tangent planes, which are complex lines for the
prepared pair.  They therefore select the same sign there.  On the
connected manifold $X$ this choice of sign cannot jump, so $B$ tames
$J_0$ everywhere.  We will choose the family of tamers after fixing its
final representative in the next step.

\paragraph{\emph{The second K\"ahler step.}}
The final pair $(A,C_1)$ is orthogonal with equal square by
\eqref{eq:endpoint-first-yau}.  Thus the structure $J_1$ already
selected by $D$ is integrable, with holomorphic volume form
$A\pm iC_1$ for the corresponding sign.  The preceding step supplies
a representative of $[B]$ taming this same $J_1$.  Even $b_1$ gives
K\"ahlerness again, and \eqref{eq:endpoint-periods} together with
\Cref{lem:endpoint-hodge-type} puts $[B]$ in type $(1,1)$.
Apply \Cref{lem:endpoint-kahler-class} with $H=[B]$ and $\nu=A^2$.
This yields a K\"ahler representative $B_1\in[B]$ such that
\begin{equation}\label{eq:endpoint-final-identities}
 AB_1=C_1B_1=0,\qquad B_1^2=A^2=C_1^2.
\end{equation}
Now apply \Cref{lem:smooth-tamers} to the family $J_t$, prescribing
the actual prepared form $B$ at $t=0$ and the newly constructed
$B_1$ at $t=1$.  It gives closed tamers $B_t\in[B]$ with these two
endpoint values.  Thus $(A,C_t,B_t)$ is a path of positive closed
triples in the original classes, joining the prepared triple to
$(A,C_1,B_1)$ itself.

Concatenate the finitely many triple paths, reparametrizing each to be
constant to all orders at its endpoints. This produces a smooth path
on $[0,1]$. Each component is closed and remains in its original
class, and positivity holds at every parameter. By
\eqref{eq:endpoint-first-yau} and
\eqref{eq:endpoint-final-identities}, if
$(\eta_1,\eta_2,\eta_3)=(A,C_1,B_1)$ and $\mu_1=A^2/2$, then
\[
 \eta_i\wedge\eta_j=2\delta_{ij}\mu_1.
\]
By \Cref{lem:complex-symplectic}, these forms are parallel and
self-dual for a hyperk\"ahler metric with volume $\mu_1$.
Finally undo the initial orthogonal change of coefficients.
Orthogonality preserves the displayed identity and restores the
original ordered classes, proving the theorem.
\end{proof}
